\scanpage{246}

\[
 \tau_{N',N}(M)\hookrightarrow M
\]

Soit $M\in\operatorname{Ob}\mathcal M_N$, $\deleted{\unsure{N}}P\in\operatorname{Ob}\mathcal M_{N'}$,

\diagram{\begin{tikzcd}[ampersand replacement=\&,column sep=small]
 \operatorname{Hom}_{\mathcal M_N}(i_{N,N'}(P),M)\arrow[d,equal]
 \& \operatorname{Hom}_{\mathcal M_{N'}}(P,\tau_{N',N}(M))\arrow[l,"\sim"']\arrow[d,equal]\\
 \operatorname{Hom}_{\mathcal M}(P,M)
 \& \operatorname{Hom}(P,\tau_{N',N}(M))\arrow[l,"\sim"']
\end{tikzcd}}

\uline{Adjoint à gauche} $q_{N',N}$ \uline{de} $i_{N,N'}$

\[
 \begin{aligned}
 \operatorname{Hom}_{\mathcal M_N}(M,i_{N,N'}(\deleted{\unsure{A}}P))&=\\[-2pt]
 \shortparallel\qquad &\\[-2pt]
 \operatorname{Hom}_{\mathcal M}(M,P)&=\operatorname{Hom}_{\mathcal M}(M/\substack{\text{ss-module}\\\text{engendré par}\\\text{les }M_i,\ i<N'},P
 \end{aligned}
\]

\diagram{\(\begin{array}{cccccccccc}
 0&0\,\cdot&M_N&M_{N+1}&M_{N+2}&\ldots&M_{N'-1}&M_{N'}&M_{N'+1}&\ldots\\
 &0&0&0&0&\ldots&0&\deleted{\unsure{M}}P_{N'}&\deleted{\unsure{M}}P_{N'+1}&\ldots
\end{array}\)}

\[
 q_{N',N}(M)=M/\substack{\text{ss-module de }M\text{ engendré}\\\text{par les }M_i,\ i<N'}
\]

\uline{NB} $q_{N',N}$ et $\tau_{N',N}$ sont des ``foncteurs de localisation'' (car adjoints d'un foncteur pl. fid.) ; \deleted{et} $q_{N',N}$ est exact à droite (en fait, \inserted{$q_{N',N},i_{N,N'},\tau_{N',N}$ commutent aux $\underrightarrow{\lim}$ quelc.}) mais en général pas à gauche sans doute -- on ne s'attend pas qu'il ait encore un adj. à g. Par contre $\tau_{N',N}$ commute aux $\underleftarrow{\lim}$ \uline{et} aux $\underrightarrow{\lim}$ (à fortiori il est exact). On s'attend à en trouver un foncteur adjoint à droite.

% archive sheet number: 245
